\section{The difference body determines the Hessian}
\label{sec:geometry}

The geometric mechanism has two parts. The objective-decreasing half of
a Dikin ellipsoid lies in the current sublevel set. Conversely, centrality
places the entire sublevel set in a bounded dilation of that ellipsoid.
Passing to differences removes the center and makes comparison between
barriers possible. The same argument tolerates a nonzero centrality
residual, measured at the point whose gap is used.

The containment ingredients are classical. The fixed-point symmetric-chord
comparison of \citet[Proposition 2.3.2(iii)]{nesterov1994} already compares
Hessians of different barriers at one interior point. Here equal objective
gap allows the compared points themselves to differ. Their use to relate sublevel
geometry and Hessian eigenvalues is also established: for primal--dual
conic sublevels, \citet[Fact 5.2 and Remark 5.1]{xiongfreund2024} bound
diameters and conic radii using central-path ellipsoids, and use this
geometry to analyze rescaling. Here we work with the tangent primal
Hessian and a prescribed primal gap. The additional conclusions proved
below are a difference-body comparison, a Loewner comparison between
different barriers at equal gap, and a matching diameter law for every
fixed barrier. The inverse-square upper edge and objective-gap
parameterization are not claimed as new; see \citet{pena2002,renegar1996}.

\subsection{Containment with an explicit centrality residual}

For $x\in P^\circ$ and $\mu>0$, define the tangent Newton decrement
\begin{equation}
 \rho_F(x;\mu)=\norm{\nabla F(x)+c_{\V}/\mu}_x^*.
 \label{eq:centrality-residual}
\end{equation}
This is the decrement of $F+\inner{c}{\cdot}/\mu$, whose Hessian is
$H_F(x)$. We call $x$ approximately central when this residual is at most
some $\rho<1$. Exact centrality is the case $\rho=0$. Define
\begin{equation}
 C(\nu,\rho)=\frac{\nu+2\sqrt\nu+\rho}{1-\rho},
 \qquad C_\nu=C(\nu,0)=\nu+2\sqrt\nu.
 \label{eq:containment-constant}
\end{equation}

\begin{lemma}[One-sided containment]
\label{lem:approx-containment}
Suppose $x\in P^\circ$, $y\in P$, and
\begin{equation}
 \inner{\nabla F(x)}{y-x}\geq-\rho\norm{y-x}_x,
 \qquad 0\leq\rho<1.
 \label{eq:one-sided-gradient}
\end{equation}
Then $\norm{y-x}_x\leq C(\nu,\rho)$. Consequently, if
$\rho_F(x;\mu)\leq\rho<1$ for some $\mu>0$ and $g=g(x)$, then
\begin{equation}
 L(g)\subset x+C(\nu,\rho)E_F(x).
 \label{eq:sublevel-containment}
\end{equation}
\end{lemma}

\begin{proof}
The assertion is immediate if $y=x$. Otherwise set
$R=\norm{y-x}_x$, $h=(y-x)/R$, and $\varphi(t)=F(x+th)$
for $0\leq t<R$. The segment is interior up to its endpoint, even when
$y$ is on the boundary. Now $\varphi''(0)=1$ and
$\varphi'(0)\geq-\rho$. The other direction of
\eqref{eq:line-hessian} gives
\begin{equation}
 \varphi''(t)\geq(1+t)^{-2},\qquad
 \varphi'(t)\geq\frac{t}{1+t}-\rho
             =\frac{(1-\rho)t-\rho}{1+t}.
 \label{eq:line-gradient-lower}
\end{equation}
Set $t_0=(\sqrt\nu+\rho)/(1-\rho)$. If $R\leq t_0$, the claimed
bound follows because $t_0<C(\nu,\rho)$. If $R>t_0$, semiboundedness
at $x+t_0h$, with endpoint $y$, gives
\[
 \varphi'(t_0)(R-t_0)\leq\nu,
 \qquad \varphi'(t_0)\geq\frac{\sqrt\nu}{1+t_0}>0.
\]
Therefore
\[
 R\leq t_0+\sqrt\nu(1+t_0)
   =\frac{\nu+2\sqrt\nu+\rho}{1-\rho}.
\]
All estimates use derivatives only at interior points, so no derivative
at $y$ is required.

For the consequence, let $e=\nabla F(x)+c_{\V}/\mu$. For $y\in L(g)$,
$\inner{c}{y-x}\leq0$, and hence
\[
 \inner{\nabla F(x)}{y-x}
 =\inner{e}{y-x}-\frac{\inner{c}{y-x}}{\mu}
 \geq-\norm{e}_x^*\norm{y-x}_x.
\]
Apply the first assertion.
\end{proof}

At $\rho=0$ this is the usual asymmetric containment with constant
$\nu+2\sqrt\nu$; this constant is also used by
\citet[Fact 5.2]{xiongfreund2024}, citing
\citet[Theorem 5.3.8]{nesterov2018}. The proof above supplies the stated
extension for every residual below one. It does not require finding an
exact center near $x$, and its sublevel is evaluated at the observed gap
$g(x)$.

\begin{theorem}[Difference-body and equal-gap comparisons]
\label{thm:difference-body}
Let $x\in P^\circ$ have gap $g$ and satisfy
$\rho_F(x;\mu)\leq\rho<1$ for some $\mu>0$. With
$C=C(\nu_F,\rho)$ and $E=E_F(x)$,
\begin{equation}
 E\subset K_g\subset2C E.
 \label{eq:difference-body-sandwich}
\end{equation}
In particular, suppose barriers $F,G$ have points $x,y$ of the same gap
$g$, with residual bounds $\rho_F,\rho_G<1$, respectively. Write
$H_F=H_F(x)$, $H_G=H_G(y)$, and
$C_F=C(\nu_F,\rho_F)$, $C_G=C(\nu_G,\rho_G)$. Then, on the fixed
tangent space,
\begin{equation}
 \frac{1}{4C_F^2}H_F\preceq H_G\preceq4C_G^2 H_F.
 \label{eq:equal-gap-loewner}
\end{equation}
The ordered eigenvalues obey the same inequalities. The condition
numbers satisfy
\begin{equation}
 \frac{\kappa_F(x)}{16C_F^2C_G^2}
 \leq\kappa_G(y)\leq16C_F^2C_G^2\kappa_F(x).
 \label{eq:equal-gap-kappa}
\end{equation}
\end{theorem}

\begin{proof}
For $h\in E$, choose $\sigma\in\{-1,1\}$ so that
$\inner{c}{\sigma h}\leq0$. Closed Dikin containment gives
$x+\sigma h\in L(g)$, while $x\in L(g)$. Thus $\sigma h\in K_g$.
The difference body is symmetric, so $h\in K_g$. This inclusion requires
no centrality. The reverse inclusion follows from
\eqref{eq:sublevel-containment}: the difference of two vectors of local
norm at most $C$ has local norm at most $2C$.

At equal gap, $E_F(x)\subset K_g\subset2C_G E_G(y)$. Applying this
to an arbitrary nonzero vector normalized in the $H_F$ norm proves
$H_G\preceq4C_G^2H_F$. Interchanging $F,G$ proves the other inequality.
The eigenvalue inequalities follow from the variational principle, and
their ratios give \eqref{eq:equal-gap-kappa}.
\end{proof}

The center itself can move substantially when the barrier changes.
Nevertheless, its Hessian ellipsoid has the same shape up to the displayed
factors, because both ellipsoids approximate the same set of chords.
This compares more than condition numbers: every quadratic form and
every ordered eigenvalue are controlled. Exact centers are covered on
their common attained gap interval from \cref{prop:gap-parametrization}.

\subsection{A matching conditioning law for every barrier}

The difference-body comparison already characterizes conditioning by an
aspect ratio without an additional instance constant. To express it using
objective accuracy, fix a closed relative ball
\begin{equation}
 z+r\BV\subset P,\qquad r>0,
 \label{eq:interior-ball}
\end{equation}
where $\BV$ is the Euclidean unit ball in $\V$. Such a ball exists because
$P^\circ$ is nonempty. Its radius and the objective range $\Delta$ measure
the fixed instance, independently of the barrier.

\begin{theorem}[Aspect ratio, spectral edges, and the diameter law]
\label{thm:diameter-law}
Under the hypotheses of \cref{thm:difference-body}, let
\[
 r_g=\max\{s\geq0:s\BV\subset K_g\}.
\]
Then $r_g>0$ and
\begin{equation}
 \left(\frac{D(g)}{2C r_g}\right)^2
 \leq\kappa_F(x)\leq
 \left(\frac{2C D(g)}{r_g}\right)^2.
 \label{eq:aspect-law}
\end{equation}
More explicitly,
\begin{align}
 \ell(x)^{-2}&\leq\lambda_{\min}(H_F(x))
          \leq C^2\ell(x)^{-2},\label{eq:minimum-rigidity}\\
 \frac{\norm{c_{\V}}^2}{g^2}
 &\leq\lambda_{\max}(H_F(x))
          \leq\left(\frac{C\Delta}{rg}\right)^2,\label{eq:maximum-rigidity}\\
 \left(\frac{\ell(x)\norm{c_{\V}}}{Cg}\right)^2
 &\leq\kappa_F(x)
          \leq\left(\frac{C\Delta\ell(x)}{rg}\right)^2.
 \label{eq:chord-law}
\end{align}
Consequently,
\begin{equation}
 \frac{\norm{c_{\V}}^2}{4C^2}\left(\frac{D(g)}{g}\right)^2
 \leq\kappa_F(x)\leq
 \frac{C^2\Delta^2}{r^2}\left(\frac{D(g)}{g}\right)^2.
 \label{eq:diameter-law}
\end{equation}
The lower bounds for $\lambda_{\min}$ and $\lambda_{\max}$ hold at
every interior point, without any centrality condition.
\end{theorem}

\begin{proof}
The Euclidean circumradius of $K_g$ is $D(g)$. An ellipsoid with Hessian
$H$ has circumradius $\lambda_{\min}(H)^{-1/2}$ and inradius
$\lambda_{\max}(H)^{-1/2}$. Apply both facts to
\eqref{eq:difference-body-sandwich} and divide the resulting inequalities
to obtain \eqref{eq:aspect-law}.

For any Euclidean unit $h\in\V$, the vector
$h/\sqrt{\inner{H_F(x)h}{h}}$ is in $E_F(x)$. Choosing its
objective-decreasing sign, as in the proof of
\cref{thm:difference-body}, shows that its Euclidean length is at most
$\ell(x)$. Thus every unit-direction Rayleigh quotient is at least
$\ell(x)^{-2}$. Conversely, choose $y\in L(g)$ with
$\norm{y-x}=\ell(x)>0$. By \eqref{eq:sublevel-containment},
the Rayleigh quotient in direction $(y-x)/\ell(x)$ is at most
$C^2/\ell(x)^2$. This proves \eqref{eq:minimum-rigidity}.

The lower bound for $\lambda_{\max}$ follows from
\eqref{eq:objective-support}, since
\[
 g^2\geq\inner{H_F(x)^{-1}c_{\V}}{c_{\V}}
      \geq\frac{\norm{c_{\V}}^2}{\lambda_{\max}(H_F(x))}.
\]
For the upper bound choose $x^*\in\Phi$. Because $0<g<\Delta$,
convexity and linearity give the homothetic inclusion
\begin{equation}
 x^*+\frac{g}{\Delta}(P-x^*)\subset L(g),
 \qquad \frac{2rg}{\Delta}\BV\subset K_g.
 \label{eq:homothetic-ball}
\end{equation}
The second inclusion follows by taking differences of the homothetic
copy of the ball in \eqref{eq:interior-ball}. Together with
$K_g\subset2C E_F(x)$ it implies, for each Euclidean unit $h$,
$\norm{(2rg/\Delta)h}_x\leq2C$. This is the upper bound in
\eqref{eq:maximum-rigidity}. Taking ratios proves
\eqref{eq:chord-law}, and \eqref{eq:chord-length} gives
\eqref{eq:diameter-law}.
\end{proof}

For geometric interpretation, the inradius $r_g$ is the minimum width of
$L(g)$, where width means the distance between parallel supporting
hyperplanes. Indeed, the support function of $K_g$ in a unit direction
$h$ is
\[
 \max_{u\in L(g)}\inner{h}{u}
 -\min_{v\in L(g)}\inner{h}{v}.
\]
A centered ball is contained in a closed convex body exactly when its
support function is no larger in every direction. Taking the minimum
over $h$ proves the assertion. Since $x\in L(g)$ has gap $g$ and
$\Phi\subset L(g)$, the objective-direction width is exactly
$g/\norm{c_{\V}}$. Thus
\begin{equation}
 \frac{2rg}{\Delta}\leq r_g\leq\frac{g}{\norm{c_{\V}}}.
 \label{eq:inradius-gap}
\end{equation}
Conditioning is therefore the squared aspect ratio of the sublevel
difference body, within barrier-parameter factors; the matching
$D(g)/g$ law results from the two bounds on its smallest width.

\begin{corollary}[Uniformity and accuracy rates]
\label{cor:uniform-rates}
For every fixed barrier on the fixed compact instance,
\begin{equation}
 \widehat\kappa_F(g)=\Theta\!\left((D(g)/g)^2\right)
 \qquad(g\downarrow0).
 \label{eq:all-barrier-rate}
\end{equation}
The same two-sided estimates hold uniformly for barriers that may depend
on $g$, provided their parameters satisfy $\nu_F\leq\bar\nu$ and the
points in question have residuals at most $\bar\rho<1$. For exact centers,
all such barriers attain every gap in
\begin{equation}
 0<g<\frac{\Delta}{1+C_{\bar\nu}}.
 \label{eq:uniform-gap-interval}
\end{equation}
If $D(g)=\Theta(g^\theta)$ with $0\leq\theta\leq1$, every fixed
barrier has $\widehat\kappa_F(g)=\Theta(g^{2\theta-2})$, and equivalently
$\kappa_F(x_F(\mu))=\Theta(\mu^{2\theta-2})$.
\end{corollary}

\begin{proof}
The first and the residual-uniform assertions follow directly from
\eqref{eq:diameter-law} with $C\leq C(\bar\nu,\bar\rho)$.
At the analytic center, $\nabla F=0$, so
\cref{lem:approx-containment} with $\rho=0$ contains all of $P$ in
$x_F^{\rm ac}+C_{\nu_F}E_F(x_F^{\rm ac})$. If $y$ maximizes the
objective on $P$, then \eqref{eq:objective-support} gives
\[
 \Delta-a_F
 =\inner{c}{y-x_F^{\rm ac}}
 \leq C_{\nu_F}\norm{c_{\V}}_{x_F^{\rm ac}}^*
 \leq C_{\nu_F}a_F.
\]
Thus $a_F\geq\Delta/(1+C_{\nu_F})$, proving
\eqref{eq:uniform-gap-interval}. The exponent follows by substitution
and \cref{prop:gap-parametrization}.
\end{proof}

The constants are explicit and separate three sources of dependence:
the barrier parameter, the allowed residual, and the fixed geometry
$(r,\Delta,\norm{c_{\V}})$. Uniformity over a family of problem
instances additionally requires control of these geometric constants.
No uniform conclusion is asserted when the barrier parameters diverge
or the residual bounds approach one. In particular, the law is an
accuracy-rate statement, not a dimension-independent runtime estimate.

\subsection{Canonical constants and compact-sublevel localization}

For canonical barriers one can sometimes improve the upper constant using
dual slack information. This is a specialization, not an additional
assumption needed for the general law. Consider either
$P=\{x:Ax=b,\ x\geq0\}$ with $F=-\sum_i\log x_i$, or
$P=\{X:AX=b,\ X_j\succeq0\ (j=1,\ldots,m)\}$ with
$F=-\sum_j\log\det X_j$. Assume the respective strictly positive or
positive-definite relative interior is nonempty, and use the Euclidean
or product Frobenius metric. Let $\nu_{\rm can}$ denote the number of
scalar coordinates or the sum of matrix-block sizes. At a central point,
the canonical dual slacks satisfy
$s_i=\mu/x_i$ or $S_j=\mu X_j^{-1}$. Hence if, on a tail
$0<\mu\leq\bar\mu$, their largest coordinate or operator norm is at
most $M$, then, with $g=g_{\rm can}(\mu)$,
\begin{equation}
 \lambda_{\max}(\widehat H_{\rm can}(g))\leq M^2\mu^{-2},\qquad
 \widehat\kappa_{\rm can}(g)
 \leq\left(\frac{M\nu_{\rm can}\ell(\widehat x_{\rm can}(g))}{g}\right)^2.
 \label{eq:canonical-constant}
\end{equation}
For the LP this follows by bounding
$\sum_i h_i^2/x_i^2$; for the SDP, bound
$\sum_j\norm{X_j^{-1/2}U_jX_j^{-1/2}}_F^2$
by $\mu^{-2}M^2\sum_j\norm{U_j}_F^2$. Use
$g\leq\nu_{\rm can}\mu$ and \eqref{eq:minimum-rigidity} to obtain
the condition-number bound.

Such a finite $M$ follows directly from any fixed strictly feasible point
$z$ or $Z$. Let $\sigma>0$ be its smallest coordinate or its smallest
block eigenvalue. Centrality supplies a dual multiplier $y_\mu$ with
$A^*y_\mu+s_\mu=c$ (using $S_\mu$ in the matrix case), where $A^*$
is the adjoint of the equality-constraint map. Canonical
complementarity gives
$\inner{b}{y_\mu}=\inner{c}{x_\mu}-\nu_{\rm can}\mu$.
Consequently,
\[
 \inner{z}{s_\mu}
 =\inner{c}{z}-\inner{b}{y_\mu}
 \leq g(z)+\nu_{\rm can}\bar\mu,
 \qquad
 M=\frac{g(z)+\nu_{\rm can}\bar\mu}{\sigma}
 \quad\text{is a valid bound},
\]
with traces and block notation understood for the SDP. Positive slacks
make the coordinate or operator-norm bound immediate. Thus this
specialized constant requires neither nondegeneracy nor strict
complementarity. In particular, the canonical barrier and every other
fixed barrier have the same conditioning rate; the equal-gap comparison
\eqref{eq:equal-gap-loewner} also gives the converse comparison of their
full spectra.

Global boundedness is convenient for central-path existence and for
choosing $\Delta$. The pointwise theorem only needs bounded lower
sublevels, which is useful for standard-form instances with unbounded
feasible regions.

\begin{proposition}[Localization to a compact sublevel]
\label{prop:localization}
Let $P$ instead be closed and convex, with nonempty relative interior,
and let the nonconstant linear objective have a finite attained optimum.
Suppose $L(\bar g)$ is compact for some $\bar g>0$. Let $F$ satisfy
\cref{def:barrier} on $P^\circ$. There is a ball
$z+r\BV\subset L(\bar g)$ with $r>0$. For every point $x\in P^\circ$
with $0<g(x)\leq\bar g$ and $\rho_F(x;\mu)\leq\rho<1$, all pointwise
conclusions of \cref{thm:difference-body,thm:diameter-law} hold with
$\Delta$ replaced by $\bar g$ in their upper bounds. The spectral-width
bounds of \cref{sec:widths} below hold as well. This assertion does not
require existence of a global central path on the unbounded set.
\end{proposition}

\begin{proof}
Take $x^*\in\Phi$ and $z_0\in P^\circ$. A sufficiently short positive
segment from $x^*$ towards $z_0$ gives $z\in P^\circ$ with
$g(z)<\bar g$. A small relative ball about $z$ lies in $P$ and below
the level $\bar g$. Dikin containment and semiboundedness are local
line arguments and remain valid. The set $L(g(x))$ is compact, so
diameters and chord maxima are attained. The proofs above now apply,
using
\[
 x^*+\frac{g}{\bar g}(L(\bar g)-x^*)\subset L(g),\qquad
 \frac{2rg}{\bar g}\BV\subset K_g
\]
in place of \eqref{eq:homothetic-ball}. All later width proofs use only
these containments and bounded objective-decreasing rays in $L(g)$.
\end{proof}
